\section{Tubular neighbourhoods of Pfaffian hypersurfaces}\label{sec:tubular}
Let $M$ be an $m$-dimensional compact Riemannian submanifold of $\R^n$ with the metric induced from the Euclidean metric, and denote by $c = n - m$ the codimension of $M$ in $\R^n$.
The $\varepsilon$-tubular neighbourhood of $M$ in $\R^n$, $T(M, \varepsilon)$, is the set of points in $\R^n$ that can be joined perpendicularly to $M$ by a segment of length $\le \varepsilon$. Since $M$ is compact, $T(M, \varepsilon)$ is equal to the closed $\varepsilon$-neighbourhood of $M$ (this need not hold if $M$ is not compact).
In this section we derive a bound on the probability that a random point lies in $T(M, \varepsilon)$ in the case where $M$ is defined by Pfaffian functions. The bound is derived by combining
a variant of Weyl's tube formula with degree bounds on the generalized Gau\ss\ map, following the approach of~\cite{lotz2015volume}.

\subsection{Generalized Gauss map}
 For $p\in M$, let $N_{p} M = \{ v \in \R^n \colon v \perp T_{p} M \}$ denote the normal space to $M$ at $p$ and  $S(NM)=\{(p,v)\in M\times S^{n-1}\colon v \perp T_{p} M\}$ the normal sphere bundle of $M$. If $M$ is a compact Riemannian manifold, then so is $S(NM)$, and the dimension of $S(NM)$ is $n-1$. 

\begin{definition}
The map $\gamma\colon S(NM) \to S^{n-1}$, $(p, v) \mapsto v$ is called the generalized Gauss map.
\end{definition}

Given a map between compact manifolds $f\colon M \to N$ of the same dimension, it is known that for a regular value $v$, the set $f^{-1}(v)$ is always finite~\cite[\S 1]{milnor1997topology}, and we define the maximum degree of $f$ by
\begin{equation*}
\md f = \max_{v \in \reg f} |f^{-1} (v)|,
\end{equation*}
where $\reg f$ is the set of regular values of $f$. 

Let $\cE^n_{j}$ be the set of dimension $j$ affine subspaces of $\R^n$. For almost all $L\in \cE^n_{c+i}$, the intersection $M\cap L$ is either empty or transverse, meaning that $\codim(T_{p}M\cap T_{p}L)=\codim T_{p}M+\codim T_{p}L$.  If $L \in \cE^n_{c+i}$ intersects $M$ transversally, then $M\cap L$ is an $i$-dimensional submanifold of $L \cong\R^{c+i}$. In this case, we can define the degree of $M$ with respect to $L$ as the maximum degree of the generalized Gauss map of $M\cap L$ as a submanifold of $L\cong \R^{c+i}$:
\begin{equation*}
\md(M, L) := \md \gamma_{M \cap L},
\end{equation*}
and the $i$-th degree of $M$ as 
\begin{equation*}
\md_i(M) := \sup_{L \in \cE^n_{c+i}} \ \md(M, L).
\end{equation*}

%\begin{figure}
%    \centering
%        \begin{tikzpicture}
%            \definecolor{lightblue}{HTML}{D9ECEA}
%            \fill [lightblue,even odd rule] (0,0) circle[radius=2.40cm] circle[radius=.60cm];
%            \draw[blue] (0,0) circle (0.60cm);
%            \draw (0,0) circle (1.5cm);
%            \draw[blue] (0,0) circle (2.40cm);
%            \node [shift={(0.95cm,0.80cm)}] {$S^1$};
%            \node [shift={(2.1cm,2.1cm)}] {$T(S^1, \varepsilon)$};
%            \draw[>=triangle 45, <->, dashed](0, 1.5cm) -- (0,2.40cm) node[midway,right] {$\varepsilon$};
%        \end{tikzpicture}
%    \caption{Tubular neighbourhood of a circle}
%    \label{fig:my_label}
%\end{figure}

A bound on the volume of $T(M,\varepsilon)$ in terms of the $i$-th degrees of $M$ was derived in~\cite[Theorem 4.3]{lotz2015volume},
based on a version of Weyl's tube formula and integral geometry. The key takeaway of this result is that the volume of $T(M, \varepsilon)$ is bounded by a polynomial in $\varepsilon$ (the original formula due to Weyl actually yields an equality, albeit with some restrictions on the size of $\varepsilon$ and different coefficients).
In what follows, $B(p,\sigma)$ denotes a closed ball of radius $\sigma$ around $p\in \R^n$.

\begin{theorem} \label{thm:tube_bound}
Let $M \subset B(p, \sigma)$ be a compact Riemannian submanifold of $\R^n$ of dimension $m$, and set $c=n-m$. Then for any $\varepsilon > 0$ we have 
\begin{equation*}
\vol T(M, \varepsilon) \le 2\omega_n \sigma^n \sum_{i=0}^m \binom{n}{c+i} \md_i(M) \left(\frac{\varepsilon}{\sigma}\right)^{c+i},
\end{equation*}
where $\omega_n = \vol B(0,1)$.
\end{theorem}

To turn this into an effective bound, we need to bound the $i$-th degrees of $M$.

\subsection{Bounding the degree of the Gauss Map}
The key to deriving an effective bound on the volume of $T(M, \varepsilon)$ is to bound the degree of the generalized Gauss map.
We are interested in the case where $M$ is a complete intersection of $c$ smooth Pfaffian functions. By this, we mean that $M$ 
can be expressed in the form
\begin{equation*}
M = \mathcal{Z}(f_1, \dots, f_c),
\end{equation*}
where the $f_i$ are Pfaffian functions of format $(\alpha, \beta, s)$, and such that the gradients $\nabla f_i$ are linearly independent at each point.

\begin{proposition}\label{prop:pfaffian-degree-bound}
  Let $M=\mathcal{Z}(f_1,\dots,f_c)$ be a Pfaffian complete intersection, where the $f_i$ are Pfaffian with format $(\alpha, \beta_i, s)$ with respect to the same Pfaffian chain and $\beta_i\le \beta$ for all $i$. Then
  \begin{align*}
    \md(M) \le{} & 2^{\frac{s(s-1)}{2}}(\alpha+\beta)^n\beta_1\cdots \beta_c \\
    &\cdot \left((n+\min\{n,s\})\alpha+n(\beta-1)+\sum_{i=1}^c(\beta_i-1)+1\right)^s.
  \end{align*}
  \end{proposition}
  
  \begin{proof}
  Let $v\in S^{n-1}$ be a regular value of the generalized Gauss map $\gamma\colon S(NM)\to S^{n-1}$. At each point $x\in M$,
  the normal sphere bundle is spanned by the normalized gradients of the $f_i$, and the number of points in the pre-image $\gamma^{-1}(v)$ is the same as the number of points in the set
  \begin{equation*}
    S:= \left\{x\in M \colon \exists \lambda_1,\dots,\lambda_c \text{ s.t. } \sum_{i=1}^c \lambda_i \nabla f_i(x) = v\right\}.
  \end{equation*}
  Note that since the gradients are linearly independent at each point, for every $x\in M$ the coefficients $\lambda_i$ are uniquely determined. The number of points in $S$ is therefore bounded by the number of solutions of the system of $n+c$ equations in $n+c$ unknowns,
  \begin{align}
    \label{eq:pfaffian-system}
  \begin{split}
    f_i(x) &= 0, \quad i\in [c]\\
    v-\sum_{i=1}^c \lambda_i \nabla f_i(x) &= 0.
  \end{split}
  \end{align}
  All functions involved are Pfaffian and we can apply Theorem~\ref{thm:khovanskii} to get a bound on the number of solutions, using the fact that the $c$ equations $f_i=0$ have degrees $\beta_i$ and the $n$ gradient equations have degree $\alpha+\beta$ by Lemma~\ref{le:pfaffalgebra}(3). The chain depends on $n$ of the $n+c$ variables (the $x$-variables only). Note that Theorem~\ref{thm:khovanskii} counts regular solutions; by Sard's theorem, for a generic regular value $v$ of $\gamma$, all solutions of~\eqref{eq:pfaffian-system} are non-degenerate. Since $|\gamma^{-1}(v)|$ is locally constant on regular values, the maximum degree is achieved at such generic values.
  \end{proof}

  \begin{remark}\label{rem:lagrange-vs-minors}
    In the proof of Proposition~\ref{prop:pfaffian-degree-bound}, one could alternatively eliminate the Lagrange multipliers $\lambda_i$ by requiring that the matrix
  \begin{equation}\label{eq:jacobian}
    J =
    \begin{bmatrix}
      \nabla f_1(x) & \cdots & \nabla f_c(x) & v
    \end{bmatrix}
  \end{equation}
  has rank at most $c$, i.e., that all $(c{+}1)\times(c{+}1)$ minors of $J$ vanish. Each such minor has degree at most $c(\alpha+\beta-1)$ in the variables $x_1,\dots,x_n$, by Lemma~\ref{le:pfaffalgebra}(3). The resulting (overdetermined) system involves only the $n$ variables $x_i$, so one can apply the connected-components bound~\eqref{eq:jones} with $D = \max\{\beta, c(\alpha+\beta-1)\}$ playing the role of $\beta$.

  While this reduces the number of variables from $n+c$ to $n$, the exponential base in~\eqref{eq:jones} becomes $\alpha + 2D - 1$, which for $c = 2$ equals $5\alpha + 4\beta - 5$; this is significantly larger than $\alpha + \beta$, the base appearing in Proposition~\ref{prop:pfaffian-degree-bound}.
  For the parameter ranges typical in neural-network applications ($c = 2$, $\alpha, \beta \le 3$), this base inflation dominates the savings from reducing the number of variables, and the Lagrange-multiplier approach gives substantially better bounds. For instance, when $\alpha = 2$, $\beta = 1$, $s = 1$, Proposition~\ref{prop:pfaffian-degree-bound} yields $\md(M) \le 3^n(2n+3)$, whereas the minor-elimination approach gives a bound of order $9^n \cdot n$. More generally, for $c = 2$ the ratio of exponential bases is $(5\alpha + 4\beta - 5)/(\alpha + \beta)$, which ranges between $4$ and $5$ for typical values of $\alpha$ and $\beta$.
\end{remark}

% \begin{lemma}\label{le:gauss-hypersurface}
% Let $V = \mathcal{Z}(f)$ be a smooth, compact hypersurface, where $f$ is a Pfaffian function with format $(\alpha, \beta, s)$. Let $\gamma\colon S(NM)\to S^{n-1}$ be the Gau\ss map of $V$. Then
% \begin{equation*}
% \md \gamma \le 2^{\frac{s(s-1)}{2} + 1} \beta (\alpha + \beta - 1)^{n-1} \big(\beta + (n-1)(\alpha + \beta - 1) - n + \min(n,s)\alpha + 1 \big)^s.
% \end{equation*}
% \end{lemma}

% \begin{proof}
% Consider the map~\footnote{This is actually the traditional definition of the Gau\ss map for an oriented hypersurface.}
% \begin{equation*}
%  \tilde{\gamma}\colon V\to S^{n-1}, \ \tilde{\gamma}(x)=\frac{\nabla f(x)}{\|\nabla f(x)\|}.
% \end{equation*}
% Any $v\in S^{n-1}$ in the image of the generalized Gau\ss map is of the form $\pm \tilde{\gamma}(x)$, so that after a suitable choice of sign for $f$, we have
% \begin{equation*}
%  |\gamma^{-1}(v)| \leq 2|\tilde{\gamma}^{-1}(v)|.
% \end{equation*} 
% %The Gauss map is given by $\gamma(x) = \frac{\nabla q}{||\nabla q||}$. We begin by bounding $\# \Tilde{g}^{-1}(v)$ for $v \in S^{n-1}$
% %a regular value. More specifically we show that 
% %$$
% %\# \Tilde{g}^{-1}(v) \le 2^{\frac{\gamma(\gamma-1)}{2}} \beta (\alpha + \beta - 1)^{n-1} \big(\beta + (n-1)(\alpha + \beta - 1) - n + \min(n,\gamma)\alpha + 1 \big)^l 
% %$$
% The argument for bounding $|\tilde{\gamma}^{-1}(v)|$ is close to the proof of the Milnor-Thom bound on the sum of Betti numbers of an algebraic variety \citep[Proposition 11.5.2]{bochnak2013real}.

% Let $\mtx{Q}$ be an orthogonal transformation such that $\mtx{Q}v=\vct{e}_n$, where $\vct{e}_n=(0,\dots,0,1)^\trans$ is the $n$th unit vector, and consider the function $g(x)=f(\mtx{Q}x)$. Then the Gauss map associated to $g$ is given by $\tilde{\gamma}_g(x)=\mtx{Q}^\trans \tilde{\gamma}(\mtx{Q}x)$, and one easily checks that 
% \begin{equation*}
%   |\{x\in \R^n \colon g(x)=0, \ \tilde{\gamma}_g(x)=\mtx{Q}v=\vct{e}_n\}|=|\{x\in \R^n \colon f(x)=0, \ \tilde{\gamma}(x)=v\}|.
% \end{equation*}
% Moreover, by Remark~\ref{re:affine} the Pfaffian structure of $g$ is the same as that of $f$. 
% We may therefore without loss of generality assume that $v = \vct{e}_n$. 
% Solutions of the equation $\Tilde{\gamma}(x) = v$ thus correspond to solutions of the system
% \begin{equation}\label{eq:paffian-hyper}
% \frac{\partial f}{\partial x_1}(x) = \dotsb = \frac{\partial f}{\partial x_{n-1}}(x) = f(x) = 0.
% \end{equation}
% Once we see that such solutions $x$ are in fact regular, which is equivalent to showing that the matrix
% \begin{equation*} 
% \mtx{J} =
% \begin{bmatrix}
%     \frac{\partial^2 f}{\partial x_1 \partial x_1} & \dots & \frac{\partial^2 f}{\partial x_1 \partial x_{n-1}} & \frac{\partial f}{\partial x_1} \\
%     \vdots                                        & \dots & \vdots                                             & \vdots                       \\ 
%     \frac{\partial^2 f}{\partial x_n \partial x_1} & \dots & \frac{\partial^2 f}{\partial x_n \partial x_{n-1}} & \frac{\partial f}{\partial x_n}
% \end{bmatrix}
% \end{equation*}
% is non-singular, we can apply Khovanskii's bound~\ref{thm:khovanskii}.

% Let $x\in \Tilde{\gamma}^{-1}(v)$. Since $\partial f(x)/\partial x_n\neq 0$, by the implicit function theorem there exists an open neighbourhood $U\subset \R^{n-1}$ and a map $h\colon U\to \R$ such that $(\vct{u},h(\vct{u}))$ parametrizes $V$ in a neighbourhood of $x$. By differentiating $f(\vct{u},h(\vct{u})) = 0$ with respect to $u_i$, we get
% \begin{equation}\label{eq:local}
%     \frac{\partial f}{\partial x_i} + \frac{\partial f}{\partial x_n} \frac{\partial h}{\partial u_i} = 0 
% \end{equation}
% in a neighbourhood of $x$. By differentiation of the above expression with respect to $u_j$, we get that 
% \begin{equation*}
%     \frac{\partial^2 f}{\partial x_j \partial x_i} = - \frac{\partial f}{\partial x_n} \frac{\partial^2 h}{\partial u_j \partial u_i}.
% \end{equation*}
% at $y$ for all $1 \le i, j \le n-1$. Using the fact that $\partial f(x)/\partial x_j=0$ for $j\in \{1,\dots,n-1\}$, the determinant
% of $\mtx{J}$ satisfies
% \begin{equation*}
%   |\det \mtx{J}| = \left|\left(\frac{\partial f(x)}{\partial x_n}\right)^n \cdot \det\left(\frac{\partial^2 h(x)}{\partial u_i\partial u_j}\right)\right|.
% \end{equation*}
% Using the expression~\eqref{eq:local}, we also get an expression of $\Tilde{\gamma}$ in the local coordinates,
% \begin{equation*}
%     \Tilde{\gamma}(u_1, \dotsc, u_{n-1}) =  \left(\frac{\partial h}{u_1}, \dotsc, \frac{\partial h}{\partial u_{n-1}}, -1 \right) \Bigg/ \sqrt{1 + \sum_{i=1}^{n-1} \Big(\frac{\partial h}{\partial u_i} \Big)^2}.
% \end{equation*}
% Differentiating and using the fact, from~\eqref{eq:local}, that $\nabla h(\vct{u}(x))=\zerovct$ at $x$, we get
% \begin{equation*}
%     \frac{\partial^2 h}{\partial u_i \partial u_j}(x) = \frac{\partial \Tilde{\gamma}_j}{\partial u_i} (x).
% \end{equation*}
% for all $1 \le i, j \le n-1$. Since $v$ is a regular value of $\tilde{\gamma}$, it follows that $\det \mtx{J}\neq 0$.

% Since $f$ is Pfaffian with format $(\alpha, \beta, s)$, it follows from Lemma~\ref{le:pfaffalgebra} that $\frac{\partial f}{\partial x_i}$ is Pfaffian with format $(\alpha, \alpha + \beta - 1, s)$, and the bound on the number of solutions of the system~\eqref{eq:paffian-hyper} follows from Khovanskii's inequality, Theorem~\ref{thm:khovanskii}. 
% \end{proof}

For $i\in [m]$, we can bound the $i$-th maximal degree as follows. Let $h\colon \R^n\to \R^n$ be an orthogonal change of variables, such that
\begin{equation*}
 L = \{x\in \R^n \colon h_{c+i+1}(x)=\cdots =h_{n}(x)=0\} 
\end{equation*}
and such that $\mathrm{mdeg}_i(M) = \mathrm{mdeg}(M;L)$. Clearly, the $i$-th maximal degree is coordinate-independent and
\begin{equation*}
  \mathrm{mdeg}_i(M) = \mathrm{mdeg}(M;L) = \mathrm{mdeg}(\tilde{M};\tilde{L}),
\end{equation*}
where $\tilde{M}=\{x\colon \tilde{f}_i(x)=f(h^{-1}(x))=0\}$ and $\tilde{L} = \{x\colon x_{c+i+1}=\cdots =x_{n}=0\}$. 
Thus $\tilde{M}$ can by seen as a submanifold of $\R^{c+i}$ and we can apply the bound from Proposition~\ref{prop:pfaffian-degree-bound} to get
\begin{align} 
  \begin{split}
  \label{eq:ideg}
  \mathrm{mdeg}_i(M) \leq &2^{\frac{s(s-1)}{2}}(\alpha+\beta)^{c+i}\beta_1\cdots \beta_c\\
  &\cdot\left((c+i+\min\{c+i,s\})\alpha+(c+i)(\beta-1)+\sum_{j=1}^c(\beta_j-1)+1\right)^s.
  \end{split}
\end{align}


\subsection{A Pfaffian tube formula}
We can now finally state and prove our bound on the volume of a tubular neighbourhood of a Pfaffian hypersurface. For convenience, this has been stated in terms of probability, 
but the reader should note that one could readily modify the proof to get a similar bound for $\vol T(V, \varepsilon)$. 

\begin{theorem} \label{thm:prob_bound}
Let $V = \mathcal{Z}(f)$ for $f$ a Pfaffian function with format $(\alpha, \beta, s)$ and $s\geq n$. Suppose that $\nabla f$ is non-vanishing on $V$ and that $V$ is bounded. Moreover, let $p \in \R^n$, $\varepsilon, \sigma > 0$ and $X$ uniformly distributed over $B(p, \sigma)$. Then
\begin{equation*}
\prob\{ d(X, V) \le \varepsilon \} \leq C_{\alpha,\beta,s,n}\left(1+(\alpha+\overline{\beta}+1)\frac{\varepsilon}{\sigma}\right)^n,
\end{equation*}
where $\overline{\beta}= \max\{\beta,2\}$ and
\begin{equation*}
C_{\alpha,\beta,s,n} =  6\cdot 2^{\frac{s(s-1)}{2}}\beta\left(n(2\alpha+\overline{\beta}-1)+\beta+1\right)^s.
\end{equation*}
\end{theorem}

\begin{proof}
The probability that $d(X,V)\leq \varepsilon$ for $X$ uniformly distributed in $B(p,\sigma)$ is given by
\begin{equation*}
  \Prob(d(X, V)\leq \varepsilon) = \frac{\vol (T(V,\varepsilon)\cap B(p,\sigma))}{\vol B(p,\sigma)}.
\end{equation*}
Let $\sigma, \varepsilon > 0$ such that $M = V \cap B(p, \sigma + \varepsilon)$ is a manifold (this is true for almost all $\sigma$), and notice that
\begin{equation*}
T(V, \varepsilon) \cap B(p, \sigma) \subset T(M \setminus \partial M, \varepsilon) \cup T(\partial M, \varepsilon),
\end{equation*}
where $\partial M=M\cap S^{n-1}(p, \varepsilon+\sigma)$ is the sphere of radius $\varepsilon+\sigma$ centered at $p$.
It is therefore enough to bound the two tubular neighbourhoods on the right-hand side.

Since $M \setminus \partial M \subset V$ and $\dim M \setminus \partial M = \dim V$, we get $\md_i(M\setminus \partial M) \leq \md_i(V)$. Although $M\setminus \partial M$ is not compact, the tube volume bound of Theorem~\ref{thm:tube_bound} extends to this setting by a compact exhaustion argument. We thus get
\begin{equation*}
  \vol T(M \setminus \partial M, \varepsilon)  \leq 2 \omega_n (\sigma+\varepsilon)^n\sum_{i=0}^{n-1} \binom{n}{i+1} \cdot \mathrm{mdeg}_i(V)\cdot \left(\frac{\varepsilon}{\sigma+\varepsilon}\right)^{i+1}.
\end{equation*}
Since $\partial M$ has codimension $2$, we get, in addition,
\begin{equation*}
  T(\partial M, \varepsilon) \leq 2 \omega_n (\sigma+\varepsilon)^{n}\sum_{i=1}^{n-1} \binom{n}{i+1} \cdot \md_{i-1}(\partial M)\cdot \left(\frac{\varepsilon}{\sigma+\varepsilon}\right)^{i+1}.
\end{equation*}
Set
\begin{equation*}
  K := 2^{\frac{s(s-1)}{2}}\beta\left(n(2\alpha+\overline{\beta}-1)+\beta+1\right)^s.
\end{equation*}
Using~\eqref{eq:ideg}, we get
\begin{align*}
\md_i(V) &\leq 2^{\frac{s(s-1)}{2}}(\alpha+\beta)^{i+1}\beta\left((i+1)(2\alpha+\beta-1)+\beta\right)^s\leq K (\alpha+\overline{\beta})^{i+1},\\
\md_i(\partial M) &\leq 2^{\frac{s(s-1)}{2}+1}(\alpha+\overline{\beta})^{i+2}\beta\left((i+2)(2\alpha+\overline{\beta}-1)+\beta+1\right)^s\leq 2K (\alpha+\overline{\beta})^{i+2}.
\end{align*}
Combining these bounds, we get
\begin{align*}
\vol T(V, \varepsilon) \cap B(p, \sigma) &\le \vol T(M \setminus \partial M, \varepsilon) + \vol T(\partial M, \varepsilon) \\
& \leq 6K\omega_n(\sigma+\varepsilon)^{n} \sum_{i=0}^{n-1}\binom{n}{i+1}(\alpha+\overline{\beta})^{i+1}\left(\frac{\varepsilon}{\sigma+\varepsilon}\right)^{i+1}.
\end{align*}
Dividing both sides by $\vol B(p, \sigma) = \omega_n \sigma^n$ and setting $a=\alpha+\overline{\beta}$ and $t=\varepsilon/\sigma$, we get
\begin{equation*}
\prob\{d(X, V)\leq \varepsilon\} \leq 6K\sum_{i=0}^{n-1}\binom{n}{i+1}(at)^{i+1}\left(1+t\right)^{n-i-1}.
\end{equation*}
The substitution $j=i+1$ and the binomial theorem give
\begin{equation*}
  \sum_{j=1}^{n}\binom{n}{j}(at)^j(1+t)^{n-j} = (1+(a+1)t)^n - (1+t)^n \leq \left(1+(a+1)\frac{\varepsilon}{\sigma}\right)^n,
\end{equation*}
which completes the proof.
\end{proof}

The proof of Theorem~\ref{thm:prob_bound} actually establishes a tighter bound that will be needed in Section~\ref{sec:robustness}.

\begin{corollary}\label{cor:tight_prob_bound}
Under the hypotheses of Theorem~\ref{thm:prob_bound}, we have
\begin{equation*}
  \prob\{ d(X, V) \le \varepsilon \} \leq C_{\alpha,\beta,s,n}\left[\left(1+(\alpha+\overline{\beta}+1)\frac{\varepsilon}{\sigma}\right)^n - \left(1+\frac{\varepsilon}{\sigma}\right)^n\right].
\end{equation*}
\end{corollary}

\begin{proof}
This follows from the penultimate line of the proof of Theorem~\ref{thm:prob_bound}, retaining the negative term in the binomial expansion.
\end{proof}

\begin{remark}
The bound in Corollary~\ref{cor:tight_prob_bound} is sharper than Theorem~\ref{thm:prob_bound} in two ways: it vanishes as $\varepsilon\to 0$ (as it must), and for large $\sigma/\varepsilon$ it decays as $C_{\alpha,\beta,s,n}\cdot n(\alpha+\overline{\beta})\varepsilon/\sigma$. This $O(\varepsilon/\sigma)$ behaviour is essential for deriving $O(1/t)$ tail bounds on condition numbers.
\end{remark}



